% Draft prepared with ChatGPT (OpenAI), 2026-09-29.
% SR coordinates follow danimalabares/daniel, grunbaum/grunbaum.tex,
% branch grunbaum-dani-understands-proof (read 2026-09-29).
% This note asks questions; it does not assert a completed smoothing proof.
\documentclass[11pt,a4paper]{article}
\usepackage[T1]{fontenc}
\usepackage{mathpazo}
\usepackage{amsmath,amssymb}
\usepackage[margin=24mm]{geometry}
\usepackage{microtype}
\usepackage{xcolor}
\usepackage{enumitem}
\usepackage[colorlinks=true,allcolors=blue!45!black]{hyperref}
\setlength{\parindent}{0pt}
\setlength{\parskip}{5pt}
\setlength{\emergencystretch}{1em}
\newcommand{\PP}{\mathbb P}
\newcommand{\II}{\mathcal I}
\DeclareMathOperator{\Hilb}{Hilb}
\DeclareMathOperator{\Proj}{Proj}
\pagestyle{plain}
\begin{document}
\begin{center}
{\Large\bfseries Linkage and a possible Calabi--Yau smoothing}\\[5pt]
{\normalsize Questions for Marcos Jardim}\\[3pt]
{\small Daniel \quad $\cdot$ \quad 29 September 2026}
\end{center}
\vspace{-6pt}

\section*{1. My example and the goal}
I am studying the Stanley--Reisner scheme of the Gr\"unbaum--Sreedharan
triangulated $3$-sphere on eight vertices. In the coordinates of my notes,
\[
S=\mathbb C[x_1,\ldots,x_8],\qquad X_0=\Proj(S/I_0)\subset\PP^7,
\]
where the sixteen minimal generators are
\[
\begin{aligned}
I_0=(&x_6x_7x_8,\ x_4x_6x_8,\ x_3x_7x_8,\ x_3x_5x_7,\\
     &x_3x_4x_8,\ x_2x_7x_8,\ x_2x_5x_7,\ x_2x_5x_6,\\
     &x_2x_4x_7,\ x_2x_4x_6,\ x_1x_4x_6,\ x_1x_4x_5,\\
     &x_1x_3x_8,\ x_1x_3x_6,\ x_1x_3x_5,\ x_1x_2x_5).
\end{aligned}
\]
This is a reduced union of twenty coordinate $\PP^3$'s, of degree $20$
and codimension $4$. Its coordinate ring is Gorenstein, with minimal
graded free resolution
\[
0\longrightarrow S(-8)\longrightarrow S(-5)^{16}
\longrightarrow S(-4)^{30}\longrightarrow S(-3)^{16}
\longrightarrow S\longrightarrow S/I_0\longrightarrow0.
\]

My hope is to smooth $X_0$ to a Calabi--Yau threefold. I believe I have
found a smoothing through deformation computations, but I am still
working through the proof, including flatness, algebraization and
smoothness of the general fibre. I therefore regard this as work in progress.

Kapustka suggested studying the special fibre via linkage. The goal of
this question is a \emph{geometric linkage description of the resulting
Calabi--Yau}: ideally, a construction from a simpler, recognizable variety.
A rational starting variety would be especially appealing.

\section*{2. A motivating example of Kapustka and collaborators}
In \cite[\S4.4]{CGKK}, start with a linear
$\Pi\cong\PP^3\subset\PP^7$ and a general complete intersection
$W=V(q_1,q_2,q_3)$ of three quadrics containing $\Pi$.
Choose a hyperplane $H=V(h)$ through $\Pi$. The residual threefold $D$
and a second residual threefold $X$ are obtained by
\[
\begin{aligned}
I_D&=(q_1,q_2,q_3,h):I_\Pi,
&\deg D&=8-1=7,\\
I_X&=(q_1,q_2,q_3,c):I_D,
&\deg X&=24-7=17,
\end{aligned}
\]
where $c$ is a suitable cubic containing $D$. Thus
$W\cap H=\Pi\cup D$ and $W\cap V(c)=D\cup X$.
Here $J:I=\{f:fI\subseteq J\}$ defines the residual scheme.

The authors obtain smooth arithmetically Gorenstein Calabi--Yau
threefolds $X$; Lemma~4.3 shows that the construction reaches a general
member of their deformation family. This degree-$17$ example illustrates
the kind of description I seek for my degree-$20$ problem.

\newpage
\section*{3. How does a chosen link vary in a family?}
Suppose, conditionally on the smoothing argument, that
$\mathcal X\subset\PP^7\times B$ is a flat embedded family over a smooth
pointed curve $(B,0)$, with special fibre $X_0$ and smooth general fibre.
Choose a codimension-four complete intersection $Z_0$ containing $X_0$,
and let $Y_0$ be its residual:
\[
I_{Z_0}=(F_1,F_2,F_3,F_4)\subset I_0,
\qquad I_{Y_0}=I_{Z_0}:I_0.
\]
\textbf{Question 1.}
Under what hypotheses does this chosen link extend along the given
deformation, after shrinking the base or an appropriate base change?
More precisely, can we lift $X_0\subset Z_0$ to
$\mathcal X\subset\mathcal Z\subset\PP^7\times B$, with $\mathcal Z$
a relative complete intersection of the same multidegree, such that
\[
\II_{\mathcal Y}=\II_{\mathcal Z}:\II_{\mathcal X},
\qquad
\II_{Y_b}=\II_{Z_b}:\II_{X_b}
\quad\text{for every }b\in B,
\]
and $\mathcal Y\to B$ is flat? I would like to understand the precise
conditions ensuring both flatness and compatibility of residual formation
with taking fibres.

In moduli language, fix the Hilbert polynomials and the multidegree
$\mathbf d$ of $Z$. Let $\mathcal F_{\mathbf d}$ be the locus in the
flag Hilbert scheme parametrizing such pairs $X\subset Z$. The relevant map is
\[
p:\mathcal F_{\mathbf d}\longrightarrow\Hilb(\PP^7),
\qquad (X\subset Z)\longmapsto[X].
\]
Is $p$ smooth at a suitable $(X_0\subset Z_0)$ in this example?
For instance, if the $F_i$ are cubics, does the ACM vanishing
$H^1(\PP^7,\II_{X_0}(3))=0$ supply the needed lifting of the four equations,
leaving relative liaison to control the residual family?

\section*{4. Can the special fibre lead to a useful description?}
\textbf{Question 2.}
Does it make sense to choose a particularly understandable link, or
sequence of links, on $X_0$ and carry it to the smooth fibres?
The crucial issue seems to be how to choose the link so that the
\emph{general residual fibre} still has recognizable geometry.

More concretely, suppose I specify a geometric family $\mathcal T$
of candidate residual threefolds. Consider the space of linkage data
$(X,Y,Z)$ with $Y\in\mathcal T$. How would one test whether its projection
to $\Hilb(\PP^7)$ dominates the component containing the proposed
smoothing branch? Which tangent-space or obstruction calculation would
be a useful first test? This would distinguish a construction of general
smooth members from a special-fibre construction alone.

\textbf{A connection with your questions.}
Your problem presentation on curves, sections of a common bundle and
their quotient sheaves made me wonder whether the linkage should be
encoded by a sheaf presentation that can be deformed. Your work on
components and smoothable boundary sheaves, for example \cite{AJT},
suggests this perspective. In this codimension-four example, is there
a useful sheaf or resolution datum to track for that purpose?
I would appreciate a criterion or a comparable example indicating
where to start.

\vspace{-5pt}
\begin{thebibliography}{9}
\small
\setlength{\itemsep}{2pt}
\bibitem{CGKK}
S.~Coughlan, \L.~Go\l\k{e}biowski, G.~Kapustka and M.~Kapustka,
\emph{Arithmetically Gorenstein Calabi--Yau threefolds in $\PP^7$},
Electron. Res. Announc. Math. Sci. \textbf{23} (2016), 52--68,
\S4.4 and Lemma~4.3.
\href{https://arxiv.org/abs/1609.01195}{arXiv:1609.01195}.
\bibitem{AJT}
C.~Almeida, M.~Jardim and A.~S.~Tikhomirov,
\emph{Irreducible components of the moduli space of rank 2 sheaves of
odd determinant on the projective space},
Adv. Math. \textbf{402} (2022), 108363.
\href{https://arxiv.org/abs/1903.00292}{arXiv:1903.00292}.
\end{thebibliography}
\end{document}
